\documentclass{article}
\usepackage[english]{babel}

%math symbols------------------------------------------------------
\usepackage{amsmath}
\usepackage{amsthm}
\usepackage{mathtools}
\usepackage{amsfonts}
\usepackage{wasysym}
\usepackage{bbm}
\usepackage{mathbbol} %lower case \mathbb
%-----------------------------------------------------------------

%referencing ----------------------------------
\usepackage{hyperref}
\usepackage[noabbrev,nameinlink,capitalise]{cleveref}
\usepackage{thmtools} %new math ambient - also for cross references
\usepackage{cite}
%-------------------------------------------

%macros ---------------------------------------
\usepackage{xargs}%multiple optional arguments
% -------------------------------------

% Graphics ----------------------------------------------
\usepackage{graphicx} % Required for inserting images
\usepackage{booktabs} 
\usepackage{multirow}
\usepackage{adjustbox}

% --------------------------------------------------------

\title{Improving chatbots\' performances using persistence}
\author{Barbara Giunti \and Nithisha Raghavaraju \and Bastian Rieck}
\date{ }

\begin{document}

\maketitle

\section{Introduction}\label{sec_intro}

Write after we're done with the experiments

\section{Background}\label{sec_background}

Barbara will take care of the PH part. 

As part of the evaluation, various conversational AI frameworks are utilized, categorized generally into classification and sequence generation approaches. The precise structural methodologies and implementations are detailed in the subsequent section.

\section{Models construction and setting}

The overarching goal of this work is to empirically investigate whether topological features derived from Persistent Homology can improve the performance and robustness of conversational AI models. We study this question across three distinct architectural paradigms: encoder-only classification models (encoder-encoder), sequence-to-sequence generation models (encoder-decoder), and autoregressive causal language models (decoder-decoder). For each paradigm, we compare a baseline model trained solely on text-derived embeddings against an augmented counterpart that additionally receives a fixed-size vector of Topological Data Analysis (TDA) features extracted from the input text. To ensure scientific rigor and isolate the impact of the topological signal, the core model capacity, hyperparameter settings, and optimization strategies remain strictly identical between the baseline and TDA-augmented versions.

\subsection{Encoder-encoder}

For the encoder-encoder classification setup, we leverage pre-trained language models to map incoming user utterances to dense contextualized vector embeddings. Specifically, we evaluate three robust foundational models. First, we use \emph{BERT} (Bidirectional Encoder Representations from Transformers), implemented specifically via the \texttt{distilbert-base-uncased} checkpoint. Second, we employ the \emph{Dense Passage Retriever (DPR)}. Finally, we evaluate the \emph{Universal Sentence Encoder (USE)}. These pre-trained models act as frozen feature extractors, and their textual embeddings serve as the primary input for a downstream shallow Multi-Layer Perceptron (MLP) classification block.

Formally, let $\mathbf{h} \in \mathbb{R}^{d}$ denote the pre-computed text embedding (e.g., the \texttt{[CLS]} token embedding, where $d = 768$ for BERT and DPR, and $d = 512$ for USE). The text embedding is passed through a fully connected projection block:
\[
  \mathbf{p} = \mathrm{Dropout}_{0.3}\left(\mathrm{ReLU}\left(\mathrm{LN}(\mathbf{W}_p \mathbf{h} + \mathbf{b}_p)\right)\right) \in \mathbb{R}^{128},
\]
where $\mathrm{LN}$ denotes Layer Normalization, and $\mathbf{W}_p$ is a learnable projection matrix. 

When topological features are incorporated, the TDA pipeline yields a feature vector $\mathbf{t} \in \mathbb{R}^{d_{\mathrm{TDA}}}$ (with $d_{\mathrm{TDA}} = 100$, obtained via Principal Component Analysis of persistence landscapes built upon FastText word embeddings). This vector is similarly projected:
\[
  \mathbf{q} = \mathrm{ReLU}(\mathrm{LN}(\mathbf{W}_q \mathbf{t})) \in \mathbb{R}^{128}.
\]
The text and TDA projections are then fused using a learned gating mechanism. This mechanism computes a soft convex combination of the representations:
\[
  \boldsymbol{\alpha} = \sigma\left(\mathbf{W}_g [\mathbf{p};\mathbf{q}]\right), \qquad \mathbf{f} = \boldsymbol{\alpha} \odot \mathbf{p} + (\mathbf{1} - \boldsymbol{\alpha}) \odot \mathbf{q},
\]
where $[\cdot;\cdot]$ denotes vector concatenation, $\sigma$ is the element-wise sigmoid function, and $\odot$ is the Hadamard product. The fused representation connects to the final classifier block. Crucially, the concatenated vector $[\mathbf{p};\mathbf{f}] \in \mathbb{R}^{256}$ is fed into the final classifier layer in both conditions; for the purely textual baseline, we set $\mathbf{f} = \mathbf{0}$, thus maintaining identical mathematical capacity within the classifier head.

Depending on the targeted dataset, the model acts as a multi-label or single-label classifier. For single-label tasks, the network minimizes the cross-entropy loss, and validation is assessed via accuracy. For multi-label scenarios, binary cross-entropy is utilized, and the model's performance is evaluated using comprehensive metrics: F1-score (Micro and Macro), Precision, Recall, and AUC. All encoder-encoder configurations are trained using the AdamW optimizer (learning rate $10^{-3}$, batch size 16) with a step-wise learning rate scheduler and an early stopping mechanism (patience of 3 epochs).

\subsection{Encoder-decoder}

The encoder-decoder architecture forms the classical backbone for open-ended chatbot response generation, mapping a user's input sequence to an appropriate conversational reply. We adapt four distinct sequence transduction models of varying inductive bias: Recurrent Neural Networks (LSTM and GRU) and Transformer-based models (a custom T4 Transformer and pre-trained BART).

\paragraph{Recurrent architectures (LSTM and GRU).} 
For the \emph{Long Short-Term Memory (LSTM)} model, tokenized input sequences are represented as word embeddings $\mathbf{e}_t \in \mathbb{R}^{E}$ ($E=256$). When TDA is enabled, the 10-dimensional TDA vector $\mathbf{t} \in \mathbb{R}^{10}$ is broadcast across the sequence length and channel-wise concatenated to the embeddings: $\tilde{\mathbf{e}}_t = [\mathbf{e}_t;\, \mathbf{t}] \in \mathbb{R}^{E + 10}$. This represents the sole architectural variation. The bidirectional LSTM encoder (hidden dimension $H=512$) condenses the sequence into a final hidden state that initializes the unidirectional LSTM decoder.
The \emph{Gated Recurrent Unit (GRU)} model follows an identical concatenation strategy at the embedding layer. However, its decoder employs a Bahdanau-style additive attention over the encoder's intermediate sequence states. During inference, both recurrent models utilize teacher-forcing during training and greedily decode the response. They are trained across 100 epochs using Adam ($10^{-3}$), equipped with gradient clipping and early stopping monitored on validation BLEU scores.

\paragraph{Transformer architectures (T4 and BART).}
The \emph{T4 Transformer} is a custom-built, randomly initialized sequence-to-sequence model featuring 4 encoder and decoder layers, 8 attention heads, and a fixed model dimension of $d_{\mathrm{model}} = 256$. To enforce strict ablation parity without altering the sub-layer dimensions, the TDA features $\mathbf{t} \in \mathbb{R}^{8}$ are injected via a tailored linear projection $\mathbf{W}_{\mathrm{TDA}} \in \mathbb{R}^{d_{\text{model}} \times \text{dim}(\mathbf{t})}$. This mapped representation is then broadcast and added element-wise to the embedded source tokens prior to sinusoidal positional encoding computation.
Conversely, \emph{BART} (\texttt{facebook/bart-base}) is a heavily parameterized ($\approx 140$M parameters) pre-trained denoising autoencoder. We introduce TDA features by linearly projecting the 10-dimensional topological vector to BART's native hidden dimension ($d_{\text{model}} = 768$) and additively injecting it into every position of the encoder's terminal hidden state before cross-attention computation in the decoder. This lightweight intervention adds a negligible $7{,}680$ parameters, effectively preserving parameter equivalence. BART fine-tuning relies on AdamW (learning rate $5\times10^{-5}$) and uses BERTScore F1 as the primary validation metric under an identical early stopping policy.

\subsection{Decoder-decoder}

The decoder-decoder regime models the conversational generation objective as a purely autoregressive language modeling task, eliminating the distinct geometric separation of the encoder. The generative models are expected to implicitly condition their response generation solely on the left-to-right causal attention prefix. 

We assess four modern foundational language models of varying scale and architecture: \texttt{distilgpt2} ($\approx 82$M parameters), \texttt{gpt2-medium} ($\approx 345$M parameters), \texttt{TinyLlama-1.1B-Chat-v1.0} ($\approx 1.1$B parameters), and \texttt{Qwen1.5-0.5B-Chat} ($\approx 0.5$B parameters). To manage the computational footprint of the larger models (i.e., TinyLlama and Qwen), we employ parameter-efficient fine-tuning via Low-Rank Adaptation (LoRA). The LoRA adapters target the query (\texttt{q\_proj}) and value (\texttt{v\_proj}) self-attention matrices with a minimal rank constraint ($r=8$) and alpha scaling factor ($\alpha=16$), heavily reducing memory overhead while preserving convergence quality. All models interpret conversational context structured by specific chat templates (e.g., encapsulating the prefix within \texttt{<|user|>} and \texttt{<|assistant|>} tokens).

Because these models lack a dedicated encoder hidden dimension for arithmetic feature fusion, topological descriptors must be seamlessly serialized. During the data preprocessing workflow, the 10-dimensional PCA-reduced topological landscape array is converted into a pseudo-lexical string payload. Each scalar coefficient $v_i$ is mapped to a discrete, format-consistent string identifier (e.g., \texttt{<tda0:0.123>} \texttt{<tda1:-0.456>}). This structured topological sequence is appended directly to the user's natural language input within the prompt bounding. Thus, the autoregressive transformer discovers the cross-modal correlations between the topological prompt tokens and the lexical content via standard causal self-attention mechanisms. Evaluation across the decoder landscape leverages generation metrics including SacreBLEU and ROUGE-L.

\section{Experiments}

\subsection{Set-up}

\begin{itemize}
    \item Description of the input data (possibly subdivided into each type - enc-enc/enc-dec/dec-dec)
    \item Technical parameter of the machine used for testing (Barbara's job)
    \item Line-up: from raw data to text embeddings (and why one type and not another), work-flow of the TDA paddling, 
    outputs
\end{itemize}

\subsection{Results and discussion}

Barbara will insert here the results once she has them

\subsubsection{Enc-enc}

\begin{table}[h!]
\begin{adjustbox}{width=\columnwidth,center}
\begin{tabular}[t]{lrrrrrrrr}
\toprule
\multirow{2}{*}{BERT} & \multicolumn{2}{c}{CounselChat} & \multicolumn{2}{c}{Depression}  & \multicolumn{2}{c}{MentalHealth} & \multicolumn{2}{c}{Suicide} \\
\cmidrule(l{3pt}r{3pt}){2-3} \cmidrule(l{3pt}r{3pt}){4-5} \cmidrule(l{3pt}r{3pt}){6-7} \cmidrule(l{3pt}r{3pt}){8-9}
 & Raw & TDA & Raw & TDA & Raw & TDA & Raw & TDA \\
\midrule 
Loss   & 0.1314 & 0.1295 & 0.4704 & 0.4777 & 0.8742 & 0.8818 & 0.2200 & 0.2334 \\ 
F1-micro   & 0.3339 & 0.3272 & 0.7580 & 0.7395 & - & - & - & - \\
F1-macro   & 0.1103 & 0.0998 & 0.6570 & 0.6451 & - & - & - & - \\
Precision  & 0.7769 & 0.8220 & 0.7768 & 0.8003 & - & - & - & - \\
Recall   & 0.2126 & 0.2042 & 0.7400 & 0.6874 & - & - & - & - \\
AUC   & 0.9114 & 0.9115 & 0.8567 & 0.8564 & - & - & - & - \\
Acc   & - & - & - & - & 0.6991 & 0.6676 & 0.9375 & 0.9261 \\ 
\bottomrule
\end{tabular}
\end{adjustbox}
\caption{}
\label{tbl:enc-enc_bert}
\end{table}

\begin{table}[h!]
\begin{adjustbox}{width=\columnwidth,center}
\begin{tabular}[t]{lrrrrrrrr}
\toprule
\multirow{2}{*}{DPR} & \multicolumn{2}{c}{CounselChat} & \multicolumn{2}{c}{Depression}  & \multicolumn{2}{c}{MentalHealth} & \multicolumn{2}{c}{Suicide} \\
\cmidrule(l{3pt}r{3pt}){2-3} \cmidrule(l{3pt}r{3pt}){4-5} \cmidrule(l{3pt}r{3pt}){6-7} \cmidrule(l{3pt}r{3pt}){8-9}
 & Raw & TDA & Raw & TDA & Raw & TDA & Raw & TDA \\
\midrule 
Loss   & 0.1229 & 0.1235 & 0.4672 & 0.4712 & 0.9359 & 0.9366 & 0.2846 & 0.1929 \\
F1-micro   & 0.4142 & 0.3770 & 0.7611 & 0.7547 & - & - & - & - \\
F1-macro   & 0.1409 & 0.1098 & 0.6721 & 0.6792 & - & - & - & - \\
Precision  & 0.7791 & 0.7815 & 0.7792 & 0.7808 & - & - & - & - \\
Recall   & 0.2821 & 0.2484 & 0.7438 & 0.7303 & - & - & - & - \\
AUC   & 0.9202 & 0.9296 & 0.8603 & 0.8575 & - & - & - & - \\
Acc   & - & - & - & - & 0.6991 & 0.6648 & 0.9148 & 0.9233 \\ 
\bottomrule
\end{tabular}
\end{adjustbox}
\caption{}
\label{tbl:enc-enc_dpr}
\end{table}

\begin{table}[h!]
\begin{adjustbox}{width=\columnwidth,center}
\begin{tabular}[t]{lrrrrrrrr}
\toprule
\multirow{2}{*}{USE} & \multicolumn{2}{c}{CounselChat} & \multicolumn{2}{c}{Depression}  & \multicolumn{2}{c}{MentalHealth} & \multicolumn{2}{c}{Suicide} \\
\cmidrule(l{3pt}r{3pt}){2-3} \cmidrule(l{3pt}r{3pt}){4-5} \cmidrule(l{3pt}r{3pt}){6-7} \cmidrule(l{3pt}r{3pt}){8-9}
 & Raw & TDA & Raw & TDA & Raw & TDA & Raw & TDA \\
\midrule
Loss   & 0.1613 & 0.1623 & 0.4365 & 0.4396 & 0.8010 & 0.8261 & 0.1865 & 0.2066 \\
F1-micro   & 0.0611 & 0.1915 & 0.7760 & 0.7814 & - & - & - & - \\
F1-macro   & 0.0112 & 0.0398 & 0.6958 & 0.6976 & - & - & - & - \\
Precision  & 0.9375 & 0.7558 & 0.8019 & 0.7896 & - & - & - & - \\
Recall   & 0.0316 & 0.7558 & 0.7516 & 0.7733 & - & - & - & - \\
AUC   & 0.8378 & 0.8428 & 0.8788 & 0.8773 & - & - & - & - \\
Acc   & - & - & - & - & 0.7479 & 0.7564 & 0.9432 & 0.9432 \\ 
\bottomrule
\end{tabular}
\end{adjustbox}
\caption{}
\label{tbl:enc-enc_use}
\end{table}

\begin{table}[h!]
\begin{adjustbox}{width=\columnwidth,center}
\begin{tabular}[t]{lrrrrrrrr}
\toprule
\multirow{2}{*}{BLEU} & \multicolumn{2}{c}{CounselChat} & \multicolumn{2}{c}{MH-Chatbot}  & \multicolumn{2}{c}{MentalChat} & \multicolumn{2}{c}{NLPMentalHealth} \\
\cmidrule(l{3pt}r{3pt}){2-3} \cmidrule(l{3pt}r{3pt}){4-5} \cmidrule(l{3pt}r{3pt}){6-7} \cmidrule(l{3pt}r{3pt}){8-9}
 & Raw & TDA & Raw & TDA & Raw & TDA & Raw & TDA \\
\midrule
distilgpt2   & 0.0336 & 0.0206 & 0.0178 & 0.0112 & 0.0349 & 0.0503 & 0.0798 & 0.1233 \\
gpt2-medium  & 0.0338 & 0.0204 & 0.0243 & 0.0154 & 0.0353 & 0.0429 & 0.0000 & 0.1482 \\
TinyLlama  & 0.0372 & 0.0188 & 0.0287 & 0.0206 &  &  &  &  \\
Qwen  & 0.0231 & 0.0197 & 0.0119 & 0.0129 &  &  &  &  \\ 
\bottomrule
\end{tabular}
\end{adjustbox}
\caption{}
\label{tbl:dec-dec_bleu}
\end{table}

\begin{table}[h!]
\begin{adjustbox}{width=\columnwidth,center}
\begin{tabular}[t]{lrrrrrrrr}
\toprule
\multirow{2}{*}{ROUGE-L} & \multicolumn{2}{c}{CounselChat} & \multicolumn{2}{c}{MH-Chatbot}  & \multicolumn{2}{c}{MentalChat} & \multicolumn{2}{c}{NLPMentalHealth} \\
\cmidrule(l{3pt}r{3pt}){2-3} \cmidrule(l{3pt}r{3pt}){4-5} \cmidrule(l{3pt}r{3pt}){6-7} \cmidrule(l{3pt}r{3pt}){8-9}
 & Raw & TDA & Raw & TDA & Raw & TDA & Raw & TDA \\
\midrule
distilgpt2   & 0.1394 & 0.1191 & 0.1010 & 0.0723 & 0.1298 & 0.1314 & 0.1151 & 0.2154 \\
gpt2-medium  & 0.1432 & 0.1260 & 0.1489 & 0.1048 & 0.1388 & 0.1215 & 0.0623 & 0.2255 \\
TinyLlama  & 0.1477 & 0.1153 & 0.1309 & 0.1097 &  &  &  &  \\
Qwen  & 0.1173 & 0.1230 & 0.0937 & 0.0595 &  &  &  &  \\ 
\bottomrule
\end{tabular}
\end{adjustbox}
\caption{}
\label{tbl:dec-dec_rougel}
\end{table}


\subsubsection{Enc-dec}

\subsubsection{Dec-dec}


\bibliographystyle{plain}
\bibliography{bibliography}

\end{document}
